%% ======================================================================
%%  Chapter II, part 2: sections 6-9 and the Appendix
%% ======================================================================

%% ----------------------------------------------------------------------
\sechead{6}{Norms and traces}{NORMS AND TRACES}

\subsechead{(a)}{REVIEW OF LINEAR ALGEBRA}

Let $A$ be a ring, $E$ a free $A$-module of finite rank and let $u$ be an
endomorphism of $E$. In linear algebra one defines the \emph{trace}, the
\emph{determinant}, and the \emph{characteristic polynomial} of $u$. If a
base $(e_{i})$ of $E$ has been chosen and if $(a_{ij})$ is the matrix for
$u$ with respect to this base, then the trace, determinant, and
characteristic polynomial of $u$ are, respectively,
\begin{equation}
\operatorname{Tr}(u)=\sum_{i=1}^{n}a_{ii},\quad
\det(u)=\det(a_{ij}),\quad\text{and}\quad
\det(X\cdot I_{\Sigma}-u)=\det(X\delta_{ij}-a_{ij}).
\label{eq:II.6.1}
\end{equation}
NB. These quantities are independent of the choice of base.

The formulas \eqref{eq:II.6.1} imply:
\begin{equation}
\begin{gathered}
\operatorname{Tr}(u+u')=\operatorname{Tr}(u)+\operatorname{Tr}(u'),\qquad
\det(uu')=\det(u)\det(u'),\\
\text{and}\qquad
\det(XI_{E}-u)=X^{n}-(\operatorname{Tr}(u))X^{n-1}+\cdots+(-1)^{n}\det(u).
\end{gathered}
\label{eq:II.6.2}
\end{equation}

\subsechead{(b)}{NORMS AND TRACES IN AN EXTENSION}

Let $B$ be a ring and let $A$ be a subring of $B$ such that $B$ is a free
$A$-module of finite rank $n$ (for example, $A$ can be a field and $B$ a
finite extension of degree $n$ of $A$). For $x\in B$, multiplication $m_{x}$
by $x$ (i.e.\ $y\mapsto xy$) is an endomorphism of the $A$-module $B$.

\begin{definition}\label{def:II.6.1}
We call \emph{trace} (respectively, \emph{norm}, \emph{characteristic
polynomial}) of $x\in B$, relative to $B$ and $A$, the trace (respectively,
determinant, characteristic polynomial) of the endomorphism $m_{x}$ of
multiplication by $x$.
\end{definition}

The trace (respectively, norm) of $x$ is denoted
$\operatorname{Tr}_{B/A}(x)$ (respectively, $\mathrm{N}_{B/A}(x)$), or
$\operatorname{Tr}(x)$ (respectively, $\mathrm{N}(x)$) when no confusion is
possible. They are elements of $A$. The characteristic polynomial is a monic
polynomial with coefficients in $A$.

For $x,x'\in B$ and $a\in A$ we have $m_{x}+m_{x'}=m_{x+x'}$ and
$m_{x}\circ m_{x'}=m_{xx'}$ and $m_{ax}=am_{x}$. Furthermore, the matrix of
$m_{a}$ with respect to any base for $B$ over $A$ is the diagonal matrix all
of whose diagonal entries are $a$. From formulas \eqref{eq:II.6.1} and
\eqref{eq:II.6.2} we obtain:
\begin{equation}
\begin{gathered}
\operatorname{Tr}(x+x')=\operatorname{Tr}(x)+\operatorname{Tr}(x'),\quad
\operatorname{Tr}(ax)=a\operatorname{Tr}(x),\quad
\operatorname{Tr}(a)=n\cdot a\\
\mathrm{N}(xx')=\mathrm{N}(x)\mathrm{N}(x'),\quad \mathrm{N}(a)=a^{n},
\quad\text{and}\quad \mathrm{N}(ax)=a^{n}\mathrm{N}(x).
\end{gathered}
\label{eq:II.6.3}
\end{equation}

\begin{proposition}
\label{prop:II.6.1}
Let $K$ be a field of characteristic $0$ or a finite field, let $L$ be an
algebraic extension of degree $n$ of $K$, let $x$ be an element of $L$, and
let $x_{1},\ldots,x_{n}$ be the roots of the minimal polynomial of $x$ over
$K$ (in a suitable extension of $K$; cf.\ \hyperref[sec:2.3]{§3}, Proposition
\ref{prop:II.3.3}), each one repeated $[L:K[x]]$ times. Then
$\operatorname{Tr}_{L/K}(x)=x_{1}+\cdots+x_{n}$,
$\mathrm{N}_{L/K}(x)=x_{1}\cdots x_{n}$. The characteristic polynomial of
$x$, relative to $L$ and $K$, is $(X-x_{1})\cdots(X-x_{n})$.
\end{proposition}

Thus the characteristic polynomial is the $[L:K[x]]$th power of the minimal
polynomial of $x$ over $K$.

\begin{proof}
Let us first treat the case where $x$ is a primitive element of $L$ over $K$
(cf.\ \hyperref[sec:2.4]{§4}, corollary of Theorem \ref{thm:II.4.1}). Let $F(X)$ be the minimal
polynomial of $x$ over $K$. Then $L$ is $K$-isomorphic to $K[X]/(F(X))$
(\hyperref[sec:2.3]{§3}, formula \eqref{eq:II.3.6}), and $(1,x,\ldots,x^{n-1})$ is a base for
$L$ over $K$. Let us put
$F(X)=X^{n}+a_{n-1}X^{n-1}+\cdots+a_{0}$. The matrix of the endomorphism
$m_{x}$ with respect to this base is
\[
M=\begin{pmatrix}
0&0&\cdots&0&-a_{0}\\
1&0&\cdots&0&-a_{1}\\
0&1&\cdots&0&\vdots\\
\vdots&0&&\vdots&\vdots\\
0&0&\cdots&1&-a_{n-1}
\end{pmatrix}.
\]
The determinant of $X\cdot I_{L}-m_{x}$ is therefore the determinant of the
matrix
\[
X\cdot I_{n}-M=\begin{pmatrix}
X&0&\cdot&&0&a_{0}\\
-1&X&&&&0&a_{1}\\
0&-1&&&&0&\vdots\\
&0&&&&\cdot\\
\cdot&\cdot&&&&X&a_{n-2}\\
0&0&&&&-1&X+a_{n-1}
\end{pmatrix}
\]
Expanding this determinant as a polynomial in $X$, we obtain the
characteristic polynomial of $x$. It is clearly equal to $F(X)$, the minimal
polynomial of $x$. By \eqref{eq:II.6.2}
$\operatorname{Tr}(x)=-a_{n-1}$ and $\mathrm{N}(x)=(-1)^{n}a_{0}$. Since $x$
is primitive, $F(X)=(X-x_{1})\cdots(X-x_{n})$; equating coefficients we see
that $\operatorname{Tr}(x)=x_{1}+\cdots+x_{n}$ and
$\mathrm{N}(x)=x_{1}\cdots x_{n}$.

Consider now the general case. Put $r=[L:K[x]]$. It suffices to show that
the characteristic polynomial $P(X)$ of $x$, with respect to $L$ and $K$, is
equal to the $r$th power of the minimal polynomial of $x$ over $K$. Let
$(y_{i})_{i=1,\ldots,q}$ be a base for $K[x]$ over $K$ and let
$(z_{j})_{j=1,\ldots,r}$ be a base for $L$ over $K[x]$; then $(y_{i}z_{j})$
is a base for $L$ over $K$ and $n=qr$ (\hyperref[sec:2.3]{§3}, Proposition \ref{prop:II.3.1}).
Let $M=(a_{ih})$ be the matrix for multiplication by $x$ in $K[x]$ with
respect to the base $(y_{i})$: thus $xy_{i}=\sum_{h}a_{ih}y_{h}$. We have
then
\[
x(y_{i}z_{j})=\Bigl(\sum_{h}a_{ih}y_{h}\Bigr)z_{j}
=\sum_{h}a_{ih}(y_{h}z_{j}).
\]
Ordering lexicographically the base $(y_{i}z_{j})$ of $L$ over $K$, we see
that the matrix $M_{1}$ for multiplication by $x$ in $L$ with respect to
this base takes the form
\[
M_{1}=\begin{pmatrix}
M&0&\cdots&0\\
0&M&\cdots&0\\
\vdots&&\vdots\vdots&\vdots\\
0&0&\cdots&M
\end{pmatrix},
\]
i.e.\ $M$ occurring $r$ times as diagonal blocks in $M_{1}$. The matrix
$X\cdot I_{n}-M_{1}$ thus consists of $r$ diagonal blocks, each of the form
$X\cdot I_{q}-M$. Consequently,
$\det(X\cdot I_{n}-M_{1})=(\det(X\cdot I_{q}-M))^{r}$. The left-hand side of
the preceding equation is $P(X)$; $\det(X\cdot I_{q}-M)$ is the minimal
polynomial of $x$ over $K$, according to the first part of the proof.
\end{proof}

\begin{proposition}
\label{prop:II.6.2}
Let $A$ be an integral domain, $K$ its field of fractions, $L$ an extension
of finite degree of $K$, and $x$ an element of $L$ integral over $A$. Assume
$K$ has characteristic zero. Then the coefficients of the characteristic
polynomial $P(X)$ of $x$ relative to $L$ and $K$, in particular
$\operatorname{Tr}_{L/K}(x)$ and $\mathrm{N}_{L/K}(x)$, are integral over
$A$.
\end{proposition}

\begin{proof}
We make use of Proposition \ref{prop:II.6.1}. We have
$P(X)=(X-x_{1})\cdots(X-x_{n})$; the coefficients of $P(X)$ are thus, up to a
sign, sums of products of the $x_{i}$'s. It suffices to show that the
$x_{i}$'s are integral over $A$ (\hyperref[sec:2.1]{§1}, Corollary \ref{cor:II.1.1} of
Proposition \ref{prop:II.1.1}). But each $x_{i}$ is a conjugate of $x$ over
$K$ (\hyperref[sec:2.4]{§4}), and there is a $K$-isomorphism
$\sigma_{i}:K[x]\to K[x_{i}]$ such that $\sigma_{i}(x)=x_{i}$. Applying
$\sigma_{i}$ to an equation of integral dependence of $x$ over $A$, we
obtain an equation of integral dependence for $x_{i}$ over $A$.
\end{proof}

\begin{statement}{Corollary}
Suppose, further, that $A$ is integrally closed. Then the coefficients of
the characteristic polynomial of $x$, in particular
$\operatorname{Tr}_{L/K}(x)$ and $\mathrm{N}_{L/K}(x)$, are elements of $A$.
\end{statement}

\begin{proof}
By definition these coefficients are elements of $K$. By Proposition
\ref{prop:II.6.2} they are integral over $A$.
\end{proof}

We remark that the quantities $x+\sigma(x)$ and $x\cdot\sigma(x)$ employed in
the discussion of quadratic fields (\hyperref[sec:2.5]{§5}) are the trace and the norm of $x$.
We proved a special case of the above corollary (\hyperref[sec:2.5]{§5}, \eqref{eq:II.5.3}) in
the course of our discussion of quadratic fields.

%% ----------------------------------------------------------------------
\sechead{7}{The discriminant}{THE DISCRIMINANT}

\begin{definition}
\label{def:II.7.1}
Let $B$ be a ring and let $A$ be a subring of $B$ such that $B$ is a free
$A$-module of finite rank $n$. For $(x_{1},\ldots,x_{n})\in B^{n}$ we call
the \emph{discriminant} of the set $(x_{1},\ldots,x_{n})$ the element of $A$
defined by the relation
\begin{equation}
\mathrm{D}(x_{1},\ldots,x_{n})=\det(\operatorname{Tr}_{B/A}(x_{i}x_{j})).
\label{eq:II.7.1}
\end{equation}
\end{definition}

\begin{proposition}
\label{prop:II.7.1}
If $(y_{1},\ldots,y_{n})\in B^{n}$ is another set of elements of $B$ such
that $y_{i}=\sum_{j=1}^{n}a_{ij}x_{j}$ with $a_{ij}\in A$, then
\begin{equation}
\mathrm{D}(y_{1},\ldots,y_{n})=(\det(a_{ij}))^{2}
\mathrm{D}(x_{1},\ldots,x_{n}).
\label{eq:II.7.2}
\end{equation}
\end{proposition}

\begin{proof}
\[
\operatorname{Tr}(y_{p}y_{q})
=\operatorname{Tr}\Bigl(\sum_{i,j}a_{pi}a_{qj}x_{i}x_{j}\Bigr)
=\sum_{i,j}a_{pi}a_{qj}\operatorname{Tr}(x_{i}x_{j}).
\]
This gives the matrix equation:
$(\operatorname{Tr}(y_{p}y_{q}))=(a_{pi})(\operatorname{Tr}(x_{i}x_{j}))
\cdot{}^{t}(a_{qj})$ (where ${}^{t}M$ denotes the transpose of the matrix
$M$). To complete the proof it suffices to take determinants.
\end{proof}

Proposition \ref{prop:II.7.1} implies that the discriminants of bases for
$B$ over $A$ are associates in $A$; i.e.\ the matrix $(a_{ij})$ which
expresses one base in terms of another has an inverse with entries in $A$.
Therefore both $\det(a_{ij})$ and $\det(a_{ij})^{-1}$ are units in $A$. We
may thus formulate the following definition:

\begin{definition}
\label{def:II.7.2}
Under the hypotheses of Definition \ref{def:II.7.1} we call the principal
ideal of $A$ generated by the discriminant of any base of $B$ over $A$ the
\emph{discriminant of $B$ over $A$}. We denote it
$\mathcal{D}_{B/A}$.
\end{definition}

\begin{proposition}
\label{prop:II.7.2}
Suppose that $\mathcal{D}_{B/A}$ contains an element which is not a
zero-divisor. Then, in order that a set
$(x_{1},\ldots,x_{n})\subset B^{n}$ be a base for $B$ over $A$, it is
necessary and sufficient that $\mathrm{D}(x_{1},\ldots,x_{n})$ generate
$\mathcal{D}_{B/A}$.
\end{proposition}

\begin{proof}
Necessity has already been proved. Suppose that
$d=\mathrm{D}(x_{1},\ldots,x_{n})$ generates $\mathcal{D}_{B/A}$. Let
$(e_{1},\ldots,e_{n})$ be a base of $B$ over $A$. Put
$d'=\mathrm{D}(e_{1},\ldots,e_{n})$ and
$x_{i}=\sum_{j=1}^{n}a_{ij}e_{j}$ with $a_{ij}\in A$, $1\leq i\leq n$. Then
$d=\det(a_{ij})^{2}d'$. By hypothesis $Ad=\mathcal{D}_{B/A}=Ad'$. Therefore,
there exists $b\in A$ such that $d'=bd$. It follows that
$d(1-b\det(a_{ij})^{2})=0$. We know that $d$ is not a divisor of zero, since
otherwise every element of $Ad=\mathcal{D}_{B/A}$ would be a divisor of
zero. Thus $1-b\det(a_{ij})^{2}=0$. This means that $\det(a_{ij})$ is
invertible. The matrix $(a_{ij})$ must be invertible, too. Consequently,
$(x_{1},\ldots,x_{n})$ is a base of $B$ over $A$.
\end{proof}

\begin{proposition}
\label{prop:II.7.3}
Let $K$ be a field which is finite or of characteristic zero, let $L$ be an
extension of finite degree $n$ of $K$, and let $\sigma_{1},\ldots,\sigma_{n}$
be the $n$ distinct $K$-isomorphisms of $L$ into an algebraically closed
field $C$ containing $K$ (\hyperref[sec:2.4]{§4}, Theorem \ref{thm:II.4.1}). Then, if
$(x_{1},\ldots,x_{n})$ is a base for $L$ over $K$,
\begin{equation}
\mathrm{D}(x_{1},\ldots,x_{n})=\det(\sigma_{i}(x_{j}))^{2}\neq0.
\label{eq:II.7.3}
\end{equation}
\end{proposition}

\begin{proof}
The first equality follows from a simple calculation.
\[
\begin{aligned}
\mathrm{D}(x_{1},\ldots,x_{n})
&=\det(\operatorname{Tr}(x_{i}x_{j}))
=\det\Bigl(\sum_{k}\sigma_{k}(x_{i}x_{j})\Bigr)
=\det\Bigl(\sum_{k}\sigma_{k}(x_{i})\sigma_{k}(x_{j})\Bigr)\\
&=\det(\sigma_{k}(x_{i}))\cdot\det(\sigma_{k}(x_{j}))
=\det(\sigma_{i}(x_{j}))^{2}.
\end{aligned}
\]
It remains to show that $\det(\sigma_{i}(x_{j}))\neq0$. We look for a
contradiction. If $\det(\sigma_{i}(x_{j}))=0$, there exist
$u_{1},\ldots,u_{n}\in C$, not all zero, such that
$\sum_{i=1}^{n}u_{i}\sigma_{i}(x_{j})=0$ for all $j$. By linearity we
conclude that $\sum_{i=1}^{n}u_{i}\sigma_{i}(x)=0$ for all $x\in L$. This
contradicts the following:
\end{proof}

\begin{statement}{Lemma of Dedekind}
Let $G$ be a group, $C$ a field, and let $\sigma_{1},\ldots,\sigma_{n}$ be
distinct homomorphisms of $G$ into the multiplicative group $C^{*}$. Then the
$\sigma_{i}$'s are linearly independent over $C$ (i.e.\
$\sum u_{i}\sigma_{i}(g)=0$ for all $g\in G$ implies that all the
$u_{i}$'s are zero).
\end{statement}

\begin{proof}
If the $\sigma_{i}$'s are linearly dependent, consider a non-trivial
relation $\sum_{i}u_{i}\sigma_{i}=0$ ($u_{i}\in C$) such that the number
$q$ of the $u_{i}$'s which are non-zero is minimum. After renumbering, we may
suppose that
\begin{equation}
u_{1}\sigma_{1}(g)+\cdots+u_{q}\sigma_{q}(g)=0\quad\text{for all }g\in G.
\label{eq:II.7.4}
\end{equation}
We have $q\geq2$, since the $\sigma_{i}$'s are not zero. For $g$ and $h$
arbitrary in $G$, we see that
\begin{equation*}
u_{1}\sigma_{1}(hg)+\cdots+u_{q}\sigma_{q}(hg)
=u_{1}\sigma_{1}(h)\sigma_{1}(g)+\cdots+u_{q}\sigma_{q}(h)\sigma_{q}(g)=0.
\end{equation*}
Multiply \eqref{eq:II.7.4} by $\sigma_{1}(h)$ and subtract. It follows that
\[
u_{2}\left(\sigma_{1}(h)-\sigma_{2}(h)\right)\sigma_{2}(g)+\cdots+
u_{q}\left(\sigma_{1}(h)-\sigma_{q}(h)\right)\sigma_{q}(g)=0.
\]
As this holds for any $g\in G$ and as $q$ has been chosen as small as
possible, it follows that
$u_{2}(\sigma_{1}(h)-\sigma_{2}(h))=0$. Thus
$\sigma_{1}(h)=\sigma_{2}(h)$ for all $h\in G$, since $u_{2}\neq0$. But this
contradicts the hypothesis that the $\sigma_{i}$'s are distinct.
\end{proof}

\begin{statement}{Remark}
Under the conditions of Proposition \ref{prop:II.7.3}, the relation
$\mathrm{D}(x_{1},\ldots,x_{n})\neq0$ means that the bilinear form
$(x,y)\mapsto\operatorname{Tr}_{L/K}(xy)$ is non-degenerate, i.e.\
$\operatorname{Tr}_{L/K}(xy)=0$ for all $y\in L$ implies $x=0$. Thus the
$K$-linear mapping which attaches to each $x\in L$ the $K$-linear form
$s_{x}:y\mapsto\operatorname{Tr}_{L/K}(xy)$ is an injection of $L$ in its
dual $\operatorname{Hom}_{K}(L,K)$ (for the structure of vector space over
$K$). As $L$ and $\operatorname{Hom}_{K}(L,K)$ are of the same finite
dimension $n$ over $K$, it follows that $x\mapsto s_{x}$ is a bijection. The
existence of ``dual bases'' of a vector space and its dual implies that, for
any base $(x_{1},\ldots,x_{n})$ of $L$ over $K$, there exists a base
$(y_{1},\ldots,y_{n})$ such that
\begin{equation}
\operatorname{Tr}_{L/K}(x_{i}y_{j})=\delta_{ij}\quad(1\leq i,j\leq n).
\label{eq:II.7.5}
\end{equation}
This remark will prove useful.
\end{statement}

\begin{theorem}
\label{thm:II.7.1}
Let $A$ be an integrally closed ring, let $K$ be its field of fractions, $L$
an extension of finite degree $n$ of $K$, and $A'$ the integral closure of
$A$ in $L$. Suppose $K$ is of characteristic $0$. Then $A'$ is an
$A$-submodule of a free $A$-module of rank $n$.
\end{theorem}

\begin{proof}
Let $(x_{1},\ldots,x_{n})$ be a base of $L$ over $K$. Each $x_{i}$ is
algebraic over $K$, so, for any $i$, we have an equation of the form
$a_{n}x_{i}^{n}+a_{n-1}x_{i}^{n-1}+\cdots+a_{0}=0$ ($a_{j}\in A$ for all
$j$). We may assume $a_{n}\neq0$. Multiplying through by $a_{n}^{n-1}$, we
see that $a_{n}x_{i}$ is integral over $A$. Put $x_{i}'=a_{n}x_{i}$. Then
$(x_{1}',\ldots,x_{n}')$ is a base for $L$ over $K$ contained in $A'$.

According to the remark preceding this theorem, there is another base
$(y_{1},\ldots,y_{n})$ of $L$ over $K$ such that
$\operatorname{Tr}(x_{i}'y_{j})=\delta_{ij}$ (\eqref{eq:II.7.5}). Let
$z\in A'$. Since $(y_{1},\ldots,y_{n})$ is a base for $L$ over $K$, we may
write $z=\sum_{j=1}^{n}b_{j}y_{j}$ with $b_{j}\in K$. For any $i$ we have
$x_{i}'z\in A'$ (since $x_{i}'\in A'$). Therefore,
$\operatorname{Tr}(x_{i}'z)\in A$ (\hyperref[sec:2.6]{§6}, corollary of Proposition
\ref{prop:II.6.2}). Thus,
$\operatorname{Tr}(x_{i}'z)=\operatorname{Tr}\bigl(\sum_{j}b_{j}
x_{i}'y_{j}\bigr)=\sum_{j}b_{j}\operatorname{Tr}(x_{i}'y_{j})
=\sum_{j}b_{j}\delta_{ij}=b_{i}$. We may conclude that $b_{i}\in A$ for all
$i$, which implies that $A'$ is a submodule of the free $A$-module
$\sum_{j=1}^{n}Ay_{j}$.
\end{proof}

\begin{statement}{Corollary}
Add to the hypotheses of Theorem \ref{thm:II.7.1} the assumption that $A$ is
principal. Then $A'$ is a free $A$-module of rank $n$.
\end{statement}

\begin{proof}
A submodule of a free $A$-module is, under our additional assumption, free
(\hyperref[sec:1.5]{Chapter I, §5}, Theorem \ref{thm:I.5.1}, (b)) and of rank $\leq n$. On the
other hand we have seen in the course of the proof of Theorem
\ref{thm:II.7.1} that $A'$ contains a base of $L$ over $K$. Therefore, it is
of rank $n$.
\end{proof}

As an exercise, the reader who lacks familiarity with the remark preceding
Theorem \ref{thm:II.7.1} should look for a more computational proof of this
theorem: with the notations defined above, set
$d=\mathrm{D}(x_{1}',\ldots,x_{n}')$ and show that, if
$z=\sum_{i}c_{i}x_{i}'$ ($c_{i}\in K$) is integral over $A$, then
$dc_{i}\in A$ (calculate $\operatorname{Tr}(zx_{j}')$ and use Cramer's
rule).

\subsechead{}{AN EXAMPLE OF THE CALCULATION OF A DISCRIMINANT}

Let $K$ be a field which is finite or of characteristic zero, let
$L=K[x]$ be an extension of finite degree $n$ of $K$, and let $F(X)$ be the
minimal polynomial of $x$ over $K$. Then
\begin{equation}
\mathrm{D}(1,x,\ldots,x^{n-1})=(-1)^{\frac{1}{2}n(n-1)}
\mathrm{N}_{L/K}(F'(x))
\label{eq:II.7.6}
\end{equation}
(where $F'(X)$ denotes the derivative of $F(X)$). Denote by
$x_{1},\ldots,x_{n}$ the roots of $F(X)$ in an extension of $K$; they are
conjugates of $x$ (\hyperref[sec:2.3]{§3}, Proposition \ref{prop:II.3.3}, and \hyperref[sec:2.4]{§4}). We see that
\[
\begin{aligned}
\mathrm{D}(1,x,\ldots,x^{n-1})
&=\det(\sigma_{i}(x^{j}))^{2}\quad\text{(by Proposition \ref{prop:II.7.3})}\\
&=\det(x_{i}^{j})^{2}=\left[\prod_{i<j}(x_{i}-x_{j})\right]^{2}
\quad\text{(Vandermonde)}\\
&=c\prod_{i\neq j}(x_{i}-x_{j})\quad
\text{(where }c=(-1)^{\frac{1}{2}n(n-1)}\text{)}\\
&=c\prod_{i}\Bigl(\prod_{j\neq i}(x_{i}-x_{j})\Bigr)\\
&=c\prod_{i}F'(x_{i})=c\,\mathrm{N}_{L/K}(F'(x))
\end{aligned}
\]
(for the $F'(x_{i})$'s are the conjugates of $F'(x)$).

In particular, applying \eqref{eq:II.7.6} to the case where $F(X)$ is a
trinomial $X^{n}+aX+b$ ($a$ and $b\in K$) and putting $y=F'(x)$, we obtain
\[
y=nx^{n-1}+a=-(n-1)a-nbx^{-1}
\]
(since $x^{n}+ax+b=0$ implies that $nx^{n-1}=-na-nbx^{-1}$). We obtain from
this $x=-nb(y+(n-1)a)^{-1}$. The minimal polynomial of $y$ over $K$ is the
numerator of $b^{-1}F(-nb(Y+(n-1)a)^{-1})$; the result of the computation is
\[
(Y+(n-1)a)^{n}-na(Y+(n-1)a)^{n-1}+(-1)^{n}b^{n-1}.
\]
The norm of $y$ is $(-1)^{n}$ times the constant term of this polynomial,
i.e.
\[
n^{n}b^{n-1}+(-1)^{n-1}(n-1)^{n-1}a^{n}.
\]
Thus,
\begin{equation}
\mathrm{D}(1,x,\ldots,x^{n-1})=(-1)^{\frac{1}{2}n(n-1)}
\left(n^{n}b^{n-1}+(-1)^{n-1}(n-1)^{n-1}a^{n}\right).
\label{eq:II.7.7}
\end{equation}
For $n=2$ (respectively, $3$) we rediscover the well-known expressions
$a^{2}-4b$ (respectively, $-27b^{2}-4a^{3}$).

%% ----------------------------------------------------------------------
\sechead{8}{The terminology of number fields}{THE TERMINOLOGY OF NUMBER FIELDS}

We call any finite (and therefore algebraic) extension of $\mathbf{Q}$ an
\emph{algebraic number field} (or \emph{number field}). For a number field
$K$, the degree $[K:\mathbf{Q}]$ is called the \emph{degree} of $K$. A number
field of degree $2$ (respectively, $3$) is called a \emph{quadratic field}
(cf.\ \hyperref[sec:2.5]{§5}) (respectively, \emph{cubic field}). A number field always has
characteristic $0$.

The elements of a number field $K$ which are integral over $\mathbf{Z}$ are
called the \emph{integers of} $K$. They form a subring $A$ of $K$ (\hyperref[sec:2.1]{§1},
Corollary \ref{cor:II.1.2} of Proposition \ref{prop:II.1.1}). This ring $A$
is a free $\mathbf{Z}$-module of rank $[K:\mathbf{Q}]$ (\hyperref[sec:2.7]{§7}, corollary of
Theorem \ref{thm:II.7.1}). The discriminants of the bases of the
$\mathbf{Z}$-module $A$ differ by a unit in $\mathbf{Z}$ (\hyperref[sec:2.7]{§7}, Definition
\ref{def:II.7.2}), a unit which is even a square in $\mathbf{Z}$ (\hyperref[sec:2.7]{§7},
Proposition \ref{prop:II.7.1}). This can only be $+1$, i.e.\ the
discriminant of the $\mathbf{Z}$-module $A$ is a well-defined element of
$\mathbf{Z}$. It is called the \emph{absolute discriminant} or the
\emph{discriminant} of $K$.

We frequently, by abuse of language, attribute to $K$ notions which are
defined relative to $A$. Thus when we speak of ideals (or units) of $K$, we
mean ideals (or units) of $A$.

%% ----------------------------------------------------------------------
\sechead{9}{Cyclotomic fields}{CYCLOTOMIC FIELDS}

We call any number field generated over $\mathbf{Q}$ by roots of unity a
\emph{cyclotomic field}. Given a prime number $p$, we write $z$ for a
primitive $p$th root of unity (in $\mathbf{C}$ for example). We are going to
study the cyclotomic field $\mathbf{Q}[z]$. The number $z$ is a root of the
polynomial $X^{p}-1$. Since $z\neq1$, it is also a root of the polynomial
$(X^{p}-1)/(X-1)=X^{p-1}+X^{p-2}+\cdots+X+1$, which is called a
\emph{cyclotomic polynomial}. It is not obvious that this polynomial is
irreducible over $\mathbf{Q}$ (i.e.\ that the field $\mathbf{Q}[z]$ is of
degree $p-1$). In order to prove that this is indeed the case we need:

\begin{statement}{Eisenstein's irreducibility criterion}
Let $A$ be a principal ideal ring, $p$ a prime element of $A$, and
\[
F(X)=X^{n}+a_{n-1}X^{n-1}+\cdots+a_{1}X+a_{0}\in A[X]
\]
such that $p$ divides $a_{i}$ ($0\leq i\leq n-1$) and $p^{2}$ does not
divide $a_{0}$. Then $F(X)$ is irreducible over the field of fractions of
$A$.
\end{statement}

\begin{proof}
Suppose that $F=GH$ with $G$ and $H\in K[X]$, both $G$ and $H$ monic
polynomials. The roots of $F$ are integral over $A$. Any root of $G$ or $H$
is a root of $F$, therefore also integral over $A$. The coefficients of $G$
(resp.\ $H$) are sums of products of roots of $G$ (resp.\ $H$); they are
therefore also integral over $A$ (\hyperref[sec:2.1]{§1}, Corollary \ref{cor:II.1.1} of
Proposition \ref{prop:II.1.1}). Since $A$ is principal, it is integrally
closed (\hyperref[sec:2.2]{§2}, \hyperref[exa:II.2.2]{Example 2}). Therefore $G\in A[X]$ and $H\in A[X]$.

Now let $\overline{F}$, $\overline{G}$, and $\overline{H}$ be the images of
$F$, $G$, and $H$ in $(A/Ap)[X]$, so
$\overline{F}=\overline{G}\cdot\overline{H}$. According to the hypothesis on
the $a_{i}$'s we have $\overline{F}=X^{n}$. Since $A/Ap$ is an integral
domain, the factorization $X^{n}=\overline{G}\cdot\overline{H}$ is
necessarily of the form $X^{n}=X^{q}\cdot X^{n-q}$ (since
$\overline{G}$ and $\overline{H}$ are monic), thus
$\overline{G}=X^{q}$ and $\overline{H}=X^{n-q}$. If $G$ and $H$ are both
non-constant, then $p$ divides the constant terms of both $G$ and $H$.
Therefore $p^{2}$ divides the constant term $a_{0}$ of $F$, and this
contradicts the hypothesis. Thus, either $G$ or $H$ is constant, and $F$ is
irreducible.
\end{proof}

\begin{statement}{Example}
The polynomial $X^{3}-2X+6$ is irreducible over $\mathbf{Q}$ (take $p=2$,
$A=\mathbf{Z}$).
\end{statement}

\begin{theorem}
\label{thm:II.9.1}
For any prime number $p$ the cyclotomic polynomial
$X^{p-1}+X^{p-2}+\cdots+X+1$ is irreducible in $\mathbf{Q}[X]$.
\end{theorem}

\begin{proof}
Put $X=Y+1$. Then
\[
\begin{aligned}
X^{p-1}+X^{p-2}+\cdots+X+1
&=\frac{X^{p}-1}{X-1}=\frac{(Y+1)^{p}-1}{Y}\\
&=Y^{p-1}+\sum_{j=p-1}^{1}\binom{p}{j}Y^{j-1}=F_{1}(Y).
\end{aligned}
\]
If $F_{1}(Y)$ is irreducible, then so is the cyclotomic polynomial.
Observing that $p$ divides each of the binomial coefficients
$\binom{p}{j}$ and that $p^{2}$ does not divide the constant term
$\binom{p}{1}=p$, we conclude that $F_{1}(Y)$ is irreducible (Eisenstein's
criterion).
\end{proof}

Theorem \ref{thm:II.9.1} implies that $\mathbf{Q}[z]$ is of degree $p-1$.
Therefore $(1,z,\ldots,z^{p-2})$ is a base for $\mathbf{Q}[z]$ over
$\mathbf{Q}$. We are going to study the ring of integers of
$\mathbf{Q}[z]$ and show that it is $\mathbf{Z}[z]$.

For this purpose we need to calculate some traces and norms (we write
$\operatorname{Tr}(x)$ and $\mathrm{N}(x)$ in place of
$\operatorname{Tr}_{\mathbf{Q}(z)/\mathbf{Q}}(x)$ and
$\mathrm{N}_{\mathbf{Q}(z)/\mathbf{Q}}(x)$). Let us note that the conjugates
of $z$ over $\mathbf{Q}$ are the $z^{j}$'s ($j=1,\ldots,p-1$) (Theorem
\ref{thm:II.9.1}).

The irreducibility of the cyclotomic polynomial implies immediately:
\begin{equation}
\operatorname{Tr}(z)=-1\quad\text{and}\quad\operatorname{Tr}(1)=p-1.
\label{eq:II.9.1}
\end{equation}
Therefore, $\operatorname{Tr}(z^{j})=-1$ for $j=1,\ldots,p-1$, and thus
\begin{equation}
\operatorname{Tr}(1-z)=\operatorname{Tr}(1-z^{2})=\cdots=
\operatorname{Tr}(1-z^{p-1})=p.
\label{eq:II.9.2}
\end{equation}
On the other hand, the calculation done in Theorem \ref{thm:II.9.1} shows
that $\mathrm{N}(z-1)=(-1)^{p-1}p$, from which it follows that
$\mathrm{N}(1-z)=p$. As the norm of $(1-z)$ is the product of the conjugates
of $1-z$, we have
\begin{equation}
p=(1-z)\bigl(1-z^{2}\bigr)\cdots\bigl(1-z^{p-1}\bigr).
\label{eq:II.9.3}
\end{equation}

Let us write $A$ for the ring of integers in $\mathbf{Q}[z]$. Evidently $A$
contains $z$ and its powers. We are going to show that
\begin{equation}
A\cdot(1-z)\cap\mathbf{Z}=p\cdot\mathbf{Z}.
\label{eq:II.9.4}
\end{equation}
We know that $p\in A\cdot(1-z)$ (\eqref{eq:II.9.3}). Thus,
$A(1-z)\cap\mathbf{Z}\supset p\cdot\mathbf{Z}$. Since $p\cdot\mathbf{Z}$ is a
maximal ideal of $\mathbf{Z}$, the relation
$A\cdot(1-z)\cap\mathbf{Z}\neq p\cdot\mathbf{Z}$ implies
$A(1-z)\cap\mathbf{Z}=\mathbf{Z}$, i.e.\ that $1-z$ is a unit in $A$. But in
this case the conjugates $(1-z^{j})$ of $1-z$ must also be units; $p$ must
be a unit in $A\cap\mathbf{Z}$ by \eqref{eq:II.9.3}; and thus $p^{-1}$ must
belong to $\mathbf{Z}$, which is absurd (\hyperref[sec:2.2]{§2}, \hyperref[exa:II.2.2]{Example 2}).

Let us show that, for any $y\in A$,
\begin{equation}
\operatorname{Tr}(y(1-z))\in p\cdot\mathbf{Z}.
\label{eq:II.9.5}
\end{equation}
Each conjugate $y_{j}(1-z^{j})$ of $y(1-z)$ is a multiple (in $A$) of
$1-z^{j}$, which is itself a multiple of $1-z$, since
$1-z^{j}=(1-z)(1+z+\cdots+z^{j-1})$. Since the trace is the sum of the
conjugates, we have $\operatorname{Tr}(y(1-z))\in A\cdot(1-z)$.
\eqref{eq:II.9.5} now follows immediately from \eqref{eq:II.9.4}, for the
trace of an integer belongs to $\mathbf{Z}$ (\hyperref[sec:2.6]{§6}, corollary of Proposition
\ref{prop:II.6.2}).

Now we are ready to determine the ring of integers of $\mathbf{Q}[z]$.

\begin{theorem}
\label{thm:II.9.2}
Let $p$ be a prime number and $z$ a primitive $p$th root of unity in
$\mathbf{C}$. Then the ring $\Lambda$ of integers of the cyclotomic field
$\mathbf{Q}[z]$ is $\mathbf{Z}[z]$, and $(1,z,\ldots,z^{p-2})$ is a base of
the $\mathbf{Z}$-module $\Lambda$.
\end{theorem}

\begin{proof}
Let $x=a_{0}+a_{1}z+\cdots+a_{p-2}z^{p-2}$ ($a_{i}\in\mathbf{Q}$) be an
element of $\Lambda$. Then
\[
x(1-z)=a_{0}(1-z)+a_{1}(z-z^{2})+\cdots+a_{p-2}(z^{p-2}-z^{p-1}).
\]
Taking traces and making use of \eqref{eq:II.9.1} and \eqref{eq:II.9.2}, we
obtain $\operatorname{Tr}(x(1-z))=a_{0}\operatorname{Tr}(1-z)=a_{0}p$. By
\eqref{eq:II.9.5}, $pa_{0}\in p\mathbf{Z}$, so $a_{0}\in\mathbf{Z}$. Since
$z^{-1}=z^{p-1}$, $z^{-1}\in\Lambda$, therefore,
\[
(x-a_{0})z^{-1}=a_{1}+a_{2}z+\cdots+a_{p-2}z^{p-3}\in\Lambda.
\]
By the same argument as before, $a_{1}\in\mathbf{Z}$. Applying the same
argument successively, we conclude that each $a_{i}\in\mathbf{Z}$, so
$\Lambda=\mathbf{Z}[z]$.
\end{proof}

\begin{statement}{Remark}
The results of this section easily extend to the case of cyclotomic fields
$\mathbf{Q}[t]$, where $t$ is a primitive $p^{r}$th root of unity ($p$
prime). Such a field is of degree $p^{r-1}(p-1)$, and its ring of integers is
$\mathbf{Z}[t]$. The minimal polynomial of $t$ over $\mathbf{Q}$ is
\[
X^{p^{r-1}(p-1)}+X^{p^{r-1}(p-2)}+\cdots+X^{p^{r-1}}+1
=\frac{X^{p^{r}}-1}{X^{p^{r-1}}-1}.
\]
\end{statement}

%% ----------------------------------------------------------------------
\annexhead{\texorpdfstring{The field $\mathbf{C}$ of complex numbers is algebraically closed}{The field C of complex numbers is algebraically closed}}%
          {THE FIELD $\mathbf{C}$ OF COMPLEX NUMBERS IS ALGEBRAICALLY CLOSED}

Let $K$ be a field and consider the following statements:
\begin{alist}
\item Any polynomial of positive degree over $K$ is a product of
polynomials of degree one (linear polynomials).
\item Any polynomial of positive degree has a root in $K$.
\end{alist}
Clearly (a) implies (b). Conversely, if (b) is true, if $P(X)$ is a
polynomial of degree $d\geq1$ over $K$, and if $a\in K$ is a root of $P(X)$,
then $P(X)$ is a multiple of $X-a$, and induction on the degree $d$ of
$P(X)$ shows that (a) is true. A field $K$ which satisfies the equivalent
conditions (a) and (b) is called \emph{algebraically closed}.

We are going to show, by a method essentially due to Lagrange, that
$\mathbf{C}$ ($=\mathbf{R}[i]$, $i^{2}=-1$) is algebraically closed. We
shall make use of only the following facts:
\begin{enumerate}[label=\arabic*.]
\item\label{fact:II.A.1} Any polynomial of odd degree over $\mathbf{R}$ has a root in
$\mathbf{R}$; this is a special case of Weierstrass's theorem on
intermediate values.
\item\label{fact:II.A.2} Any quadratic polynomial over $\mathbf{C}$ has its roots in
$\mathbf{C}$. The elementary ``formula'' for a root of $aX^{2}+bX+c=0$
reduces the question to that of showing that any $z=a+ib\in\mathbf{C}$
($a,b\in\mathbf{R}$) has a square root in $\mathbf{C}$. But
$(x+iy)^{2}=a+ib$ ($x,y\in\mathbf{R}$) is equivalent to $x^{2}-y^{2}=a$ and
$2xy=b$; it follows that $a^{2}+b^{2}=(x^{2}+y^{2})^{2}$ and
$x^{2}+y^{2}=\sqrt{a^{2}+b^{2}}$. Clearly, one may find $x^{2}$ and $y^{2}$,
hence $x$ and $y$, satisfying these equations.
\item\label{fact:II.A.3} Given a non-constant polynomial $P(X)\in K[X]$, there exists an
extension $K'$ of $K$ such that $P(X)$ factors into a product of linear
polynomials in $K'[X]$. There was an easy proof of this fact in Proposition
\ref{prop:II.3.3} of \hyperref[sec:2.3]{§3} (a proof almost independent of other results in this
chapter; it suffices to know that, if $F(X)$ is irreducible,
$K[X]/(F(X))$ is a field. Then use an induction).
\item\label{fact:II.A.4} The relations between coefficients and roots of a polynomial.
\item\label{fact:II.A.5} The fact that a symmetric polynomial
$G(X_{1},\ldots,X_{r})\in K[X_{1},\ldots,X_{n}]$ is a polynomial in the
elementary symmetric functions
$\sum X_{i},\sum X_{i}X_{j},\ldots,X_{1}\cdots X_{n}$ of the $X_{i}$'s.
\end{enumerate}

\begin{statement}{Theorem}
The field $\mathbf{C}$ of complex numbers is algebraically closed.
\end{statement}

\begin{proof}
We shall prove property (b), that any non-constant polynomial has a root in
$\mathbf{C}$. Observing that
$F(X)=P(X)\overline{P}(X)$ ($\overline{P}$ the polynomial whose
coefficients are the complex conjugates of the coefficients of $P$) has
roots in $\mathbf{C}$ if and only if $P$ does, we see that we need consider
only polynomials with real coefficients.

Now we write the degree of $F(X)$ ($\in\mathbf{R}[X]$) in the form
$d=2^{n}q$ where $q$ is odd. We argue by induction on the exponent $n$ of
$2$. For $n=0$, $d$ is odd and $F(X)$ has a root in $\mathbf{R}$ (cf.\ \hyperref[fact:II.A.1]{(1)}).
Suppose $n\geq1$. By \hyperref[fact:II.A.3]{(3)} there exists an extension $K'$ of $\mathbf{C}$ and
$x_{1},\ldots,x_{d}\in K'$ such that
$F(X)=\prod_{i=1}^{d}(X-x_{i})$ (assuming, without loss of generality, that
$F(X)$ is monic).

Let $c$ be an arbitrary element of $\mathbf{R}$ and consider the elements
$y_{ij}=x_{i}+x_{j}+cx_{i}x_{j}$ of $K'$ ($i\leq j$). The cardinality of
$(y_{ij})$, $1\leq i\leq j\leq d$, is
$\frac{1}{2}d(d+1)=2^{n-1}q(d+1)$, where $q(d+1)$ is odd. The polynomial
$G(X)=\prod_{i\leq j}(X-y_{ij})$ has as coefficients symmetric polynomials
in the $x_{i}$'s with real coefficients. By \hyperref[fact:II.A.5]{(5)} the coefficients are
polynomials in the elementary symmetric functions of the $x_{i}$'s; these
polynomials themselves have real coefficients. Therefore, the coefficients
of $G(X)$ are real (by \hyperref[fact:II.A.4]{(4)}). As its degree is of the form $2^{n-1}\times$
(odd), the induction hypothesis implies that it has a root
$z_{c}\in\mathbf{C}$. One of the $y_{ij}$'s, say
$y_{i(c)j(c)}=x_{i(c)}+x_{j(c)}+cx_{i(c)}x_{j(c)}$, is therefore equal to
$z_{c}$.

Now, since $\mathbf{R}$ is infinite and since the set of pairs $(i,j)$
($i\leq j$) is finite, there exist two distinct real numbers $c$ and $c'$
such that $i(c)=i(c')$ and $j(c)=j(c')$. Denote these indices by $r$ and
$s$, respectively. Then
$x_{r}+x_{s}+cx_{r}x_{s}=z_{c}\in\mathbf{C}$ and
$x_{r}+x_{s}+c'x_{r}x_{s}=z_{c'}\in\mathbf{C}$. Taking linear combinations,
we may conclude that $x_{r}+x_{s}\in\mathbf{C}$ and
$x_{r}x_{s}\in\mathbf{C}$. Therefore, by \hyperref[fact:II.A.4]{(4)}, $x_{r}$ and $x_{s}$ are roots
of a quadratic equation with coefficients in $\mathbf{C}$. We may conclude
that $x_{r},x_{s}\in\mathbf{C}$ (by \hyperref[fact:II.A.2]{(2)}). Thus $F(X)$ has a root in
$\mathbf{C}$ and the theorem is proved.
\end{proof}
